\documentclass[../../../include/open-logic-section]{subfiles}
%READERNOTE{SOL6-A281: एकच द्विमान निरूपण-पद्धत वापरून एकास-एक फलनाच्या सिद्धतेतील अंक-एकमेवपणाची अट पूर्ण केली आहे. मूळ प्रतिमेबाहेरील संख्येचे उदाहरण [0,1] वर प्रतिमेतच मिळते; खरे प्रतिउदाहरण 5/6 दिले आहे. हा आधुनिक द्विमान युक्तिवाद मूळ दशमान पत्रातील शब्दशः सिद्धता म्हणून मांडलेला नाही.}

\begin{document}

\olfileid{his}{set}{cantorplane}
\olsection{रेषा आणि प्रतल: कँटर यांचा निष्कर्ष}

श्रेडर--बर्नस्टाइन प्रमेयाच्या सिद्धतेशी संबंधित काही परिस्थितींचा संबंध \olref[his][set][pathology]{sec} मध्ये चर्चिलेल्या इतिहासाशी आहे. आठवा: १८७७ मध्ये कँटर यांनी चौरसावर जितके बिंदू आहेत तितकेच त्याच्या एका बाजूवरही आहेत, हे सिद्ध केले. त्यांच्या सुरुवातीच्या प्रयत्नातील अंक एकाआड एक लिहिण्याची कल्पना येथे द्विमान रूपात मांडू. मूळ ऐतिहासिक प्रयत्न दशमान विस्तारांवर आधारलेला होता; खालील द्विमान मांडणी ही त्याचे आधुनिक रूपांतर आहे (\cite[पृ.~201--202]{Gouvea2011}).

$\unitline$ ही एकक रेषा, म्हणजे $[0,1]$ या बिंदूंचा संच, असू द्या. $\unitsquare$ हा एकक चौरस, म्हणजे $\unitline \times \unitline$ या बिंदूंचा संच, असू द्या. या संज्ञांत कँटर यांनी $\cardeq{\unitline}{\unitsquare}$ सिद्ध केले. त्यांनी डेडेकिंड यांना पत्रातून कळवलेल्या कल्पनेचे द्विमान रूपांतर पुढे दिले आहे.

\begin{thm}\ollabel{thm:cantorplane}
$\cardeq{\unitline}{\unitsquare}$
\end{thm}

\begin{proof}[सिद्धता: पहिला भाग.]
$a,b\in\unitline$ निश्चित करा. द्विमान विस्तारासाठी पुढील
एकच पद्धत ठरवू: $0=0.000\ldots$; दोन निरूपणे असलेल्या
धन संख्येसाठी शेवटी अनंत $1$ असलेले निरूपण वापरू.
विशेषतः $\nicefrac12=0.0111\ldots$ आणि $1=0.111\ldots$.
इतर संख्येचे द्विमान निरूपण एकमेव असते. या निवडीने
प्रत्येक धन संख्येच्या विस्तारात अनंत $1$ येतात. असे लिहा:
\begin{align*}
a &= 0.a_1a_2a_3a_4\dots\\
b &= 0.b_1b_2b_3b_4\dots
\intertext{आता पुढीलप्रमाणे दिलेले $f\colon\unitsquare\to\unitline$ हे फलन विचारात घ्या:}
f(a,b)&=0.a_1b_1a_2b_2a_3b_3a_4b_4\dots
\end{align*}
परिणामी विस्तारात अनंत $1$ असतात, जोपर्यंत $a=b=0$
नाहीत; त्या प्रकरणात तो पूर्णपणे शून्यांचा असतो.
म्हणून परिणामी विस्तारही वर ठरवलेल्या पद्धतीतीलच आहे.
त्याच्या अंकांचे एकमेवपण वापरल्यास, $f(a,b)=f(c,d)$
असताना प्रत्येक $n\in\Nat$ साठी $a_n=c_n$ आणि
$b_n=d_n$ मिळते. म्हणून $a=c$ आणि $b=d$;
म्हणजे $f$ हे !!a{injection} आहे.
\end{proof}

मात्र हे फलन !!a{surjection} नाही. उदाहरणार्थ,
द्विमान संख्या $z=0.1101010101\ldots=\nicefrac56$ घ्या.
तिचा द्विमान विस्तार सांत नाही, म्हणून त्या विस्ताराची दुसरी
निवड उपलब्ध नाही. $f(a,b)=z$ असते, तर तिचे विषम आणि
सम स्थानांवरील अंक वेगळे केल्याने अनुक्रमे
\begin{align*}
a&=0.100000\ldots=\nicefrac12,\\
b&=0.111111\ldots=1
\end{align*}
असे विस्तार मिळाले असते. पण $a$ चा हा विस्तार आपण
ठरवलेल्या पद्धतीला मान्य नाही: $\nicefrac12$ साठी
$0.011111\ldots$ वापरायचा आहे. त्यामुळे $z$ साठी
कोणतीही जोडी $\tuple{a,b}$ मिळत नाही.
उदाहरणार्थ, $0.101010\ldots$ हा मात्र
या पद्धतीत $f(1,0)$ आहे; म्हणून तो येथे प्रतिउदाहरण नाही.

अंक एकाआड एक लिहिण्याच्या ऐतिहासिक प्रयत्नात
निरूपणांची निवड महत्त्वाची आहे, हे डेडेकिंड यांनी
निदर्शनास आणले. वरच्या निश्चित द्विमान पद्धतीत $f$
हे !!a{injection} आहे, पण !!a{surjection} नाही.
म्हणून या पहिल्या भागाने फक्त
$\cardle{\unitsquare}{\unitline}$ दाखवले आहे;
$\cardeq{\unitsquare}{\unitline}$ अजून सिद्ध झालेले नाही.

कँटर यांनी डेडेकिंड यांना जवळजवळ लगेच उत्तर लिहून सिद्धता साधारणपणे पुढीलप्रमाणे पूर्ण करता येईल असे सुचवले:

\begin{proof}[सिद्धता: पूर्ण रूप.]
आपण $\cardle{\unitsquare}{\unitline}$ दाखवले आहे. पण $\unitline$ वरून $\unitsquare$ मध्ये जाणारे !!a{injection} उघडच आहे: रेषा चौरसाच्या एका बाजूवर सरळ ठेवा. म्हणून $\cardle{\unitline}{\unitsquare}$ आणि $\cardle{\unitsquare}{\unitline}$. श्रेडर--बर्नस्टाइन प्रमेयानुसार (\olref[sfr][siz][sb]{thm:schroder-bernstein}) $\cardeq{\unitline}{\unitsquare}$.
\end{proof}

मात्र कँटर यांना अर्थातच शेवटची ओळ या रूपात पूर्ण करता आली नसती, कारण श्रेडर--बर्नस्टाइन प्रमेय तेव्हा अजून सिद्ध झाले नव्हते. पुढे कँटर यांनी हा सर्वसाधारण अंदाज मांडला, तरी त्याची समाधानकारक सिद्धता १८९७ पर्यंत मिळाली नव्हती. (म्हणूनच १८७७ मध्ये पुढे कँटर यांनी \olref{thm:cantorplane} ची श्रेडर--बर्नस्टाइन प्रमेयाचा आधार न घेणारी वेगळी सिद्धता दिली.)

\end{document}
